\title{xLSTM: Extended Long Short-Term Memory}

\begin{abstract}
The foundational concepts of the Long Short-Term Memory (LSTM) architecture—namely, the constant error carousel and gating mechanisms—were first proposed in the 1990s. Since their inception, LSTMs have maintained their relevance, undergirding multiple breakthroughs in deep learning, especially serving as the basis for the original Large Language Models (LLMs). The emergence of Transformer architectures, however, fundamentally shifted the landscape, owing to their scalable and parallelizable self-attention modules, thereby eclipsing LSTMs at large scales. Here, we pose a straightforward research question: To what extent can LSTMs be advanced in the field of language modeling if they are scaled to billions of parameters, incorporate state-of-the-art techniques from recent LLM developments, and address traditional weaknesses? To this end, we first propose the adoption of exponential gating, enhanced through specialized normalization and stabilization strategies. Next, we reengineer the core LSTM memory component, leading to: (i) sLSTM, which employs scalar memory, a scalar-based update, and a novel scheme for memory mixing, and (ii) mLSTM, designed for full parallelization using a matrix-form memory paired with a covariance-based update mechanism. By embedding these LSTM variants within residual block frameworks, we form xLSTM blocks, which can be stacked residually to build xLSTM models. The introduction of exponential gating and refined memory architectures strengthens xLSTM’s abilities, enabling these models to rival and even surpass leading Transformers and State Space Models in both predictive performance and scalability.\\
Code available at: \url{https://github.com/NX-AI/xlstm}
\end{abstract}

\section{Introduction}
\label{sec:intro}

The Long Short-Term Memory (LSTM) ideas~\citep{Hochreiter:91,Hochreiter:97nips,Hochreiter:97}, i.e., the constant error carousel and gating, were introduced to overcome the vanishing gradient problem of recurrent neural networks~\citep{Hochreiter:91,Hochreiter:00book}:

\begin{align}
\label{eq:lstm_idea}
  c_t \ = \ f_t \ c_{t-1} \ + \ 
          i_t \  z_t \ , \quad  
          h_t \ = \ o_t \ \psi({c_t}) \ . 
\end{align}

The constant error carousel operates through an additive modification to the cell state $c_{t-1}$ (depicted in green) using the cell inputs $z_t$, which is regulated via sigmoid gating mechanisms (shown in blue). The input gate $i_t$ alongside the forget gate $f_t$ modulate this update process, whereas the output gate $o_t$ governs the production of the memory cell’s output, corresponding to the hidden state $h_t$. The cell state is transformed or compressed by $\psi$, and this processed state is subsequently passed through the output gate to yield the hidden state.

LSTMs have achieved significant success across multiple fields \citep{Hochreiter:01,Hochreiter:07,Schmidhuber:15}, dominating text generation tasks up until the emergence of Transformers in 2017~\citep{Vaswani:17}. Their utility has been proven in a variety of sequence-related tasks such as text generation~\citep{Graves:13,Karpathy:15blog}, handwriting synthesis~\citep{Graves:13}, sequence-to-sequence translation~\citep{Sutskever:14nips}, program evaluation~\citep{Zaremba:14arxiv}, captioning images~\citep{Karpathy:15,Hossain:19}, generating source code~\citep{Karpathy:15blog}, rainfall-runoff predictions~\citep{Kratzert:18,Kratzert:19}, and developing hydrological flood warning models~\citep{Nearing:24}. Within reinforcement learning, LSTMs continue to represent the leading approach for sequence modeling, exemplified by their use in high-profile models such as AlphaStar for StarCraft II~\citep{Vinyals:17short}, OpenAI Five for Dota 2~\citep{OpenAI:19}, and controllers for nuclear fusion magnets~\citep{Degrave:22short}. LSTMs possess a strong capability for abstraction, effectively encoding semantic representations and maintaining them within their memory cells~\citep{Karpathy:15blog}, as demonstrated by the discovery of neurons specializing in number and syntax~\citep{Lakretz:19}, linguistic properties~\citep{Bau:19}, or sentiment~\citep{Radford:17arxiv}. Even today, LSTMs continue to underpin crucial applications~\citep{Degrave:22short,Nearing:24}, demonstrating their enduring relevance.

Although LSTMs have achieved significant progress, they exhibit three primary shortcomings: (i) Inability to update stored information. This is illustrated using the {\it Nearest Neighbor Search} task (further details in Appendix~\ref{sec:appNearestNeighborSearch}). Here, given a reference vector, a model processes a sequence to identify and return the associated value of the vector most similar to the reference at the end of the sequence. As depicted in the left panel of Figure~\ref{fig:lstmProblems}, the mean squared error indicates that LSTMs struggle to replace previously stored values with better matches encountered later in the sequence. In contrast, our newly introduced xLSTM employs exponential gating to successfully address this issue. (ii) Constrained memory resources; specifically, all data must be encoded in scalar cell states. This shortcoming is demonstrated through the {\it Rare Token Prediction} task. The right panel of Figure~\ref{fig:lstmProblems} presents token prediction perplexities on Wikitext-103~\citep{Merity:17}, grouped by the frequency of tokens. LSTMs exhibit higher perplexity for infrequent tokens, a direct consequence of their limited storage. Our xLSTM incorporates a matrix memory architecture that mitigates this deficit. (iii) A lack of computational parallelism arises due to memory interdependency—hidden-to-hidden connections link sequential steps, enforcing strict order in processing.

Such drawbacks have contributed to the widespread adoption of Transformers~\citep{Vaswani:17} in language modeling. This prompts the following question: If LSTMs are redesigned to circumvent these limitations and are scaled up to match contemporary Large Language Models, what improvements in language modeling performance might be achievable?

\section{Extended Long Short-Term Memory}
\label{sec:xLSTM}

To address the shortcomings of LSTM, Extended Long Short-Term Memory (xLSTM) incorporates two significant changes to the LSTM concept presented in Equation~\eqref{eq:lstm_idea}. Specifically, xLSTM introduces exponential gating and innovative memory architectures. These enhancements lead to the emergence of two novel variants within the LSTM framework: (i) sLSTM (refer to Section~\ref{sec:sLSTM}), characterized by the use of a scalar memory, scalar updates, and the incorporation of memory mixing, and (ii) mLSTM (see Section~\ref{sec:mLSTM}), which employs a memory represented as a matrix and utilizes a covariance (outer product) update scheme, enabling complete parallelism. Exponential gating serves to strengthen both the sLSTM and mLSTM models. In pursuit of parallel execution, the mLSTM discards memory mixing, meaning it omits recurrent connections among hidden states. Both mLSTM and sLSTM permit extension to configurations with multiple memory cells; in this scenario, sLSTM retains memory mixing across distinct cells. Additionally, the sLSTM variant may support several heads, facilitating memory mixing among cells within individual heads while precluding mixing across different heads. Introducing heads alongside exponential gating into sLSTM paves the way for a novel approach to memory mixing. Conversely, for mLSTM, configurations with multiple heads and those with multiple memory cells are functionally equivalent.

Embedding these enhanced LSTM structures within residual blocks leads to the formation of xLSTM blocks (refer to Section~\ref{sec:xLSTMblock}). When these xLSTM blocks are stacked using residual connections, they compose the building blocks of xLSTM architectures (refer to Section~\ref{sec:xLSTMarchitecure}). Figure~\ref{fig:xlstm_sketch} provides an illustration of the xLSTM architecture and its fundamental components.

\subsection{Review of the Long Short-Term Memory}
\label{sec:orgLSTM}

The foundational concept behind LSTM~\citep{Hochreiter:91,Hochreiter:97nips,Hochreiter:97} proposed a scalar memory cell that functions as the primary unit for both storage and computation, specifically designed to mitigate the vanishing gradient problem~\citep{Hochreiter:91,Hochreiter:00book} via the constant error carousel mechanism—that is, through the cell state update. This memory cell architecture integrates three distinct gates: the input gate, the output gate, and, subsequently introduced by \citet{Gers:00}, the forget gate.  
The dynamics of the LSTM memory cell at a given time point $t$ are described by the following update equations:

\begin{align}
c_t \ &= \  f_t \ c_{t-1} \ + \ i_t \ z_t &  & &\text{cell state} \\
h_t  \ &= \ o_t \ \tilde{h}_t \ , & \tilde{h}_t \ &= \ \psi \left( c_t\right)
  &\text{hidden state} \\
z_t \ &= \ \varphi \left( \tilde{z}_t \right) \ , 
  &\tilde{z}_t \ &=  \ \mathbf{w}^\top_{z} \ \mathbf{x}_t \ + \
  r_{z}  h_{t-1} \ + \  b_{z} \ \
  &\text{cell input} \\
i_t \ &= \ \sigma \left( \tilde{i}_t  \right) \ , 
  &\tilde{i}_t \ &= \ \mathbf{w}^\top_{i} \ \mathbf{x}_t \ + \
  r_{i}  \ h_{t-1} \ + \  b_{i} \ \
  &\text{input gate} \\
f_t \ &= \ \sigma \left(  \tilde{f}_t \right) \ , 
  &\tilde{f}_t \ &= \ \mathbf{w}^\top_{f} \ \mathbf{x}_t  \ + \
  r_{f}  \ h_{t-1} \ + \  b_{f} \ \
  &\text{forget gate} \\
o_t \ &= \ \sigma \left( \tilde{o}_t \right) \ , 
  &\tilde{o}_t  \ &= \ \mathbf{w}^\top_{o} \ \mathbf{x}_t \ + \
  r_{o}  \ h_{t-1} \ + \  b_{o} \ \
  &\text{output gate} 
\end{align}

The vectors $\mathbf{w}_{z}$, $\mathbf{w}_{i}$, $\mathbf{w}_{f}$, and $\mathbf{w}_{o}$ serve as the input weights linking the input $\mathbf{x}_t$ to the cell input, input gate, forget gate, and output gate, accordingly. Likewise, the parameters $r_{z}$, $r_{i}$, $r_{f}$, and $r_{o}$ represent the recurrent weights connecting the previous hidden state $h_{t-1}$ to the same respective components: cell input, input gate, forget gate, and output gate. The terms $b_{z}$, $b_{i}$, $b_{f}$, and $b_{o}$ denote the biases for each corresponding connection. The activation functions for the cell input and hidden state, $\varphi$ and $\psi$ (usually instantiated as $\tanh$), play essential roles; specifically, $\psi$ regulates the cell state, preventing it from diverging. Gate activations are carried out with the sigmoid function, $\sigma \left( x \right)= 1/(1+\exp(-x))$. In more recent versions, the memory is structured as a vector of scalar cells $\mathbf{c}_t \in \mathbb{R}^{d}$, which enables the deployment of recurrent weight matrices $\mathbf{R} \in \mathbb{R}^{d \times d}$ to interconnect and blend outputs from different memory cells~\citep{Greff:15}; further explanation is provided in Appendix~\ref{sec:apporgLSTM}. Experiments involving ablation have indicated that the full memory cell architecture is essential for proper function~\citep{Greff:15}.

\subsection{sLSTM}
\label{sec:sLSTM}

In order to enable LSTMs to reconsider their storage choices, we propose incorporating exponential gates (depicted in red), along with mechanisms for normalization and stabilization. Specifically, the input and forget gates are designed to utilize exponential activation functions. For the normalization component, a normalizer state is established, which accumulates the sum of the products obtained by multiplying the input gate with every subsequent forget gate.

The scalar sLSTM forward pass is:

\begin{align}
\label{eq:slstmforward}
c_t \ &= \  f_t \ c_{t-1} \ + \ i_t \ z_t & & 
 &\text{cell state} \\
n_t \ &= \  f_t \ n_{t-1} \ + \ i_t  & & 
 &\text{normalizer state} \\
h_t \ &= \ o_t \ \tilde{h}_t \ ,
  &\tilde{h}_t \ &= \ c_t / n_t
&\text{hidden state} \\
z_t \ &= \ \varphi \left( \tilde{z}_t \right) \ , 
  &\tilde{z}_t \ &=  \ \mathbf{w}^\top_{z} \ \mathbf{x}_t \ + \
  r_{z}  h_{t-1} \ + \  b_{z}
  &\text{cell input} \\
i_t \ &= \ \!\exp \left( \tilde{i}_t  \right) \ , 
  &\tilde{i}_t \ &= \ \mathbf{w}^\top_{i} \ \mathbf{x}_t \ + \
  r_{i}  \ h_{t-1} \ + \  b_{i}
  &\text{input gate} \\
f_t \ &= \ \sigma \left(  \tilde{f}_t \right) \ \text{OR} \
\!\exp \left(  \tilde{f}_t \right) \ ,
  &\tilde{f}_t \ &= \ \mathbf{w}^\top_{f} \ \mathbf{x}_t  \ + \
  r_{f}  \ h_{t-1} \ + \  b_{f}
  &\text{forget gate} \\
o_t \ &= \ \sigma \left( \tilde{o}_t \right) \ , 
  &\tilde{o}_t  \ &= \ \mathbf{w}^\top_{o} \ \mathbf{x}_t \ + \
  r_{o}  \ h_{t-1} \ + \  b_{o}
  &\text{output gate} 
\end{align}

We transfer the original LSTM gating techniques, i.e., input- and/or hidden-dependent gating plus bias term, to the new architectures. Exponential activation functions can lead to large values that cause overflows. Therefore, we stabilize gates with an additional state \mbox{$m_t$~\citep{Milakov:18arxiv}}:

\begin{align}
\label{eq:slstmstabil}
m_t \ &= \ \max \left( \log ( f_t ) + m_{t-1} , \log ( i_t ) \right) &\text{stabilizer state} \\
i'_t \ &= \ \!\exp \left( \log \left ( i_t \right) - m_t \right) \ = \ \!\exp \left( \tilde{i}_t - m_t \right) \ \, 
  &\text{stabil. input gate} \\
f'_t \ &= \ \!\exp \left( \log \left( f_t \right) + m_{t-1} - m_t \right) \ \, 
  &\text{stabil. forget gate}
\end{align}

As demonstrated in Appendix~\ref{sec:appsLSTM}, substituting $f_t$ with $f'_t$ and $i_t$ with $i'_t$ during the forward computation has no impact on the overall network’s output or on the gradients of the loss function with respect to its parameters.

\paragraph{New Memory Mixing.} Similar to the original LSTM, sLSTM is capable of utilizing several memory cells (refer to Appendix~\ref{sec:appsLSTM}). The presence of multiple memory cells allows for the mixing of memory through recurrent pathways $\mathbf{R}_{\mathbf{z}}$, $\mathbf{R}_{\mathbf{i}}$, $\mathbf{R}_{\mathbf{f}}$, and $\mathbf{R}_{\mathbf{o}}$ connecting the hidden state $\mathbf{h}$ to the memory cell input $\mathbf{z}$, as well as to the gating vectors $\mathbf{i}$, $\mathbf{f}$, and $\mathbf{o}$. Notably, exponential gating brings a novel dynamic to memory mixing. In this revised architecture, sLSTM permits several heads that perform memory mixing internally but do not share memory mixing across different heads. The integration of heads in sLSTM, combined with exponential gating, introduces a distinctive memory mixing scheme.

\subsection{mLSTM}
\label{sec:mLSTM}

To expand the storage capabilities of LSTMs, we transition the memory cell from a scalar representation $c \in \mathbb{R}$ to a matrix form $\mathbf{C} \in \mathbb{R}^{d \times d}$. Consequently, retrieval operations are executed through matrix multiplication. At each time step $t$, a pair of vectors is to be memorized—specifically, the key $\mathbf{k}_t \in \mathbb{R}^d$ and the value $\mathbf{v}_t \in \mathbb{R}^d$, consistent with the terminology used in Transformers. Subsequently, at a future step $t+\tau$, the stored value $\mathbf{v}_t$ should be accessed using a query vector $\mathbf{q}_{t+\tau} \in \mathbb{R}^d$. This context aligns with the framework of Bidirectional Associative Memories (BAMs) \citep{Kohonen:72,Anderson:72,Nakano:72,Anderson:77}.

The covariance update rule~\citep{Sejnowski:77,Dayan:91} for storing a key-value pair is

\begin{align}
\mathbf{C}_t \ &= \ \mathbf{C}_{t-1} \ + \ \mathbf{v}_{t} \ \mathbf{k}_t^\top \ .
\end{align}

We posit that applying layer normalization prior to projecting inputs into keys and values results in zero-mean representations. The covariance update rule achieves optimality~\citep{Dayan:91} with respect to maximizing the separability among the stored binary vectors, which corresponds to optimizing the signal-to-noise ratio. Enhanced separability is attainable when restricting retrieval operations to pairwise associations, even though this incurs quadratic computational complexity akin to the attention mechanism~\citep{Krotov:16,Krotov:17,Ramsauer:21}. Notably, the covariance update rule coincides with the methodology of Fast Weight Programmers~\citep{Schmidhuber:92ncfastweights,Schlag:21}. Subsequent work has incorporated a fixed decay factor scaling $\mathbf{C}_{t-1}$ and a constant learning rate applied to $\mathbf{v}_{t} \mathbf{k}_t ^\top$~\citep{Ba:16}. Building on these developments, we incorporate the covariance update rule within the LSTM paradigm: the forget gate represents the decay parameter, the input gate acts as the learning rate, and the output gate modulates the outputted retrieved vector.

For this matrix memory, the normalizer state is the weighted sum of key vectors, where each key vector is weighted by the input gate and all future forget gates. Again, the normalizer state keeps record of the strength of the gates. Since the dot product between query and normalizer state can be close to zero, we use the absolute value of this dot product and lower bound it by a threshold (typically 1.0) as done previously~\citep{Sun:23arxiv}. The mLSTM forward pass is:

\begin{align}
\label{eq:mlstm_recurrent_begin}
\mathbf{C}_t \ &= \  f_t \ \mathbf{C}_{t-1} \ + \ 
  i_t \ \mathbf{v}_t \ \mathbf{k}_t^\top & &  &\text{cell state} \\
\mathbf{n}_t \ &= \  f_t \ \mathbf{n}_{t-1} \ + \ 
  i_t \ \mathbf{k}_t & &  &\text{normalizer state} \\
\mathbf{h}_t  \ &= \ \mathbf{o}_t \ \odot \ \tilde{\mathbf{h}}_t
\ , \qquad \qquad 
\tilde{\mathbf{h}}_t \ = \   \mathbf{C}_t \mathbf{q}_t \ / \ 
  \max \left\{ |\mathbf{n}_t^\top \mathbf{q}_t|, 1 \right\} 
   &&&\text{hidden state} \\
\mathbf{q}_t \ &= \ \mathbf{W}_q \ \mathbf{x}_t \ + \ \mathbf{b}_q  & &  &\text{query input} \\
\mathbf{k}_t \ &= \ \frac{1}{\sqrt{d}} \mathbf{W}_k \ \mathbf{x}_t \ + \ \mathbf{b}_k  & &  &\text{key input} \\
\mathbf{v}_t \ &= \ \mathbf{W}_v \ \mathbf{x}_t \ + \ \mathbf{b}_v  & &  &\text{value input} \\
i_t \ &= \ \!\exp \!  \! \left( \tilde{i}_t  \right) \ , 
 \qquad \qquad \ \ \ \ \,
  \tilde{i}_t \ = \ \mathbf{w}^\top_{i} \ \mathbf{x}_t \ + \  b_{i} \ \ \ 
  &&&\text{input gate} \\
f_t \ &= \ \sigma \!  \left(  \tilde{f}_t \right) \ \text{OR} \ \!\exp \!  \! \left(  \tilde{f}_t \right) , \ \ \: \!
\tilde{f}_t \ = \ \mathbf{w}^\top_{f} \ \mathbf{x}_t  \ + \
  b_{f}
  &&& \text{forget gate} \\
\label{eq:mlstm_recurrent_end}
\mathbf{o}_t \ &= \ \sigma \left( \tilde{\mathbf{o}}_t \right) \ , \qquad \qquad \quad \ \ \ \,
  \tilde{\mathbf{o}}_t  \ = \ \mathbf{W}_{\mathbf{o}} \ \mathbf{x}_t \ + \
  \mathbf{b}_{\mathbf{o}}
  &&&\text{output gate} 
\end{align}

Similar to the original LSTM, mLSTM supports multiple memory cells. Within mLSTM architectures, employing several heads or several memory cells yields the same effect due to the absence of interactions among the memory cells. To ensure stability of the exponential gating mechanisms in mLSTM, the same stabilization procedures applied to sLSTM are utilized (refer to Equation~\ref{eq:slstmstabil}). Furthermore, given that mLSTM lacks memory mixing, its recurrence relations can be re-expressed in a parallelized form. Additional information is provided in Appendix~\ref{sec:appmLSTM}.

\subsection{xLSTM Architecture}
\label{sec:xLSTMarch}

\paragraph{xLSTM Blocks.}
\label{sec:xLSTMblock}
The purpose of an xLSTM block is to non-linearly encapsulate the sequence’s past within a high-dimensional representation, thereby facilitating greater separability between distinct contextual trajectories. Distinguishing these trajectories is essential for accurate sequence prediction tasks, such as forecasting the subsequent token. To justify this architectural approach, we draw upon Cover’s Theorem~\citep{Cover:65}, which posits that non-linear embedding of data into spaces of increased dimensionality enhances the likelihood of their linear separability compared to the original lower-dimensional space. Within this framework, we analyze two types of residual block structures: (i) The post up-projection residual block (akin to that used in Transformers), wherein the historical information is first non-linearly condensed within the original feature space, followed by a linear transformation to a higher-dimensional domain; a subsequent non-linear activation is then applied, after which the data is linearly projected back to its initial dimensionality. This architecture is illustrated in the left panel of Figure~\ref{fig:Backbones} and in the third column of Figure~\ref{fig:xlstm_sketch}. An expanded schematic of this configuration is provided in Figure~\ref{fig:appsLSTM_detailed} within the appendix. (ii) The pre up-projection residual block (as seen in State Space Models), which commences with a linear projection into a high-dimensional space,

compresses previous information non-linearly within a high-dimensional space, followed by a linear transformation returning it to its original space. When an xLSTM block includes an sLSTM, we employ a post up-projection module; conversely, with an mLSTM-based xLSTM block, the pre up-projection module is utilized due to the increased memory in high-dimensional representation. For further clarification, see the left panel of Figure~\ref{fig:Backbones}, the third column of Figure~\ref{fig:xlstm_sketch}, or Figure~\ref{fig:appsLSTM_detailed} in the appendix.

\paragraph{xLSTM Architecture. }
\label{sec:xLSTMarchitecure}
The xLSTM architecture is formed by stacking constituent blocks in a residual manner~\citep{Srivastava:15,He:16}. Our approach incorporates the widely adopted pre-LayerNorm~\citep{Ba:16b} residual backbone design common to state-of-the-art Large Language Models. Refer to the final column of Figure~\ref{fig:xlstm_sketch} for a visualization.

\subsection{Memory and Speed Considerations}

In contrast to Transformer architectures, xLSTM networks exhibit linear computational complexity and maintain constant memory requirements regardless of the input sequence length. Owing to the compressive nature of xLSTM’s memory system, this architecture proves advantageous for deployment in industrial environments and on edge devices.

The memory used in mLSTM does not involve learnable parameters, but its operation involves intensive computation due to reliance on a $d \times d$ matrix memory and associated $d \times d$ updates. This setup prioritizes expanded memory capacity at the cost of greater computational burden. Nonetheless, since these operations are amenable to parallelization on GPUs, their impact on actual processing time remains relatively modest.

Although mLSTM lends itself to parallel execution—akin to approaches like FlashAttention~\citep{Dao:22,Dao:23} and GLA~\citep{Yang:23arxiv}—the sLSTM variant is not similarly parallelizable as a result of memory mixing through hidden-to-hidden connections. To address this, we have engineered a high-speed CUDA implementation with register-level GPU memory optimization, resulting in execution speeds that are typically no more than double those of the mLSTM baseline.

\section{Related Work}
\label{sec:related}

\paragraph{Linear Attention.} To address the challenge of quadratic time complexity in Transformer models with respect to context length, a variety of linear attention techniques have been introduced. The Synthesizer eliminates explicit token-to-token dependencies by learning artificial attention weights~\citep{Tay:20arxiv}. The Linformer implements self-attention via a projection into a lower-dimensional space, allowing for a linear approximation~\citep{Wang:20arxiv}. The Linear Transformer redesigns the attention mechanism so that it can be computed in linear time \citep{Katharopoulos:20}. Performer leverages positive orthogonal random features to linearly approximate the softmax function used in attention~\citep{Choromanski:21}. Meanwhile, alternatives such as Structured Global Convolution (SGConv)~\citep{Li:22} and Hyena Hierarchy~\citep{Poli:23} substitute the attention operation with efficient long-range convolutional modules.

\paragraph{State Space Models.} State Space Models (SSMs) have garnered considerable attention in recent times, owing to their linear computational complexity with respect to context length and their competitive results relative to Transformer architectures. Among the earliest examples is the Structured State Space sequence model (S4)~\citep{Gu:21}. Subsequent developments include the Diagonal State Space (DSS) model~\citep{Gupta:22}, Gated State Space (GSS) models~\citep{Mehta:22}, the S5 model~\citep{Smith:22}, Bidirectional Gated SSM (BiGS)~\citep{Wang:22}, H3 model~\citep{Fu:23}, as well as Mamba~\citep{Gu:24arxiv}.

\paragraph{Recurrent Neural Networks.} Recurrent Neural Networks (RNNs) have been proposed as alternatives to Transformers and attention mechanisms, principally because their computational complexity scales linearly with sequence length. Recent advancements in RNNs employing Deep Linear Recurrent Units (LRUs) have demonstrated strong performance in language modeling tasks~\citep{Orvieto:23, De:24arxiv}. Notable architectures such as the Hierarchically Gated Linear RNN (HGRN)~\citep{Qin:23} and its successor HGRN2~\citep{Qin:24arxiv} have achieved similar success. The RWKV framework~\citep{Peng:23arxivshort,Peng:24arxivshort}, a prominent RNN-based solution for large-scale language models, has yielded results comparable to those achieved by Transformer-based methods.

\paragraph{Gating.}
The concept of gating, integral to the LSTM architecture, has been independently rediscovered and reframed within a number of contemporary models. Variants utilizing gating mechanisms include HGRN~\citep{Qin:23}, HGRN2~\citep{Qin:24arxiv}, Gated Linear Attention (GLA)~\citep{Yang:23arxiv}, Gated State Space (GSS) models~\citep{Mehta:22}, Bidirectional Gated SSM (BiGS)~\citep{Wang:22}, Moving Average Equipped Gated Attention (MEGA)~\citep{Ma:22}, RWKV~\citep{Peng:23arxivshort}, and Mamba~\citep{Gu:24arxiv}.

\paragraph{Covariance Update Rule.} To support improved memory retention, our approach incorporates a matrix-based memory within the mLSTM cell, updated via a covariance update procedure. Related strategies that utilize such update rules include Fast Weight Programmers~\citep{Schmidhuber:92ncfastweights,Schlag:21}, RWKV-5 and RWKV-6~\citep{Peng:24arxivshort}, Retention~\citep{Sun:23arxiv}, Linear Transformer~\citep{Katharopoulos:20}, as well as HGRN2~\citep{Qin:24arxiv}.

\paragraph{Most Related.} The models most similar in spirit to xLSTM are Retention~\citep{Sun:23arxiv}, RWKV~\citep{Peng:23arxivshort,Peng:24arxivshort}, and HGRN2~\citep{Qin:24arxiv}. These frameworks all utilize the ideas of matrix memory and/or gating mechanisms. Nonetheless, unlike the newly introduced sLSTM, they do not incorporate memory mixing. The inclusion of memory mixing is pivotal for tackling state tracking challenges, making LSTMs fundamentally more capable than State Space Models (SSMs) and Transformers~\citep{Merrill:24,Deletang:23} with respect to expressiveness. State tracking forms a key component in tasks such as program execution or following entities over long texts.

\paragraph{Residually Stacking Architectures.} Consistent with the design paradigm of most state-of-the-art deep learning systems, xLSTM models employ residual stacking of core modules~\citep{Srivastava:15,He:16}.

This residual stacking approach is foundational to the success of deep convolutional networks~\citep{He:16}, as well as the Transformer framework~\citep{Vaswani:17}. Transformers, in turn, underpin the most powerful Large Language Models (LLMs) such as GPT-3~\citep{Brown:20short}, ChatGPT~\citep{Schulman:22short}, GPT-4~\citep{Achiam:23gpt4short}, Megatron-LM~\citep{Shoeybi:19}, Gopher~\citep{Rae:21short}, ERNIE 3.0 Titan~\citep{Wang:21arxivshort}, GLaM~\citep{Du:21glamshort}, Chinese M6~\citep{Lin:21}, the multilingual AlexaTM 20B~\citep{Soltan:22}, OPT~\citep{Zhang:22opt}, Chinchilla~\citep{Hoffmann:22short}, BLOOM~\citep{Scao:22arxivshort}, GLM-130B~\citep{Zeng:22arxivshort}, LaMDA~\citep{Thoppilan:22short}, PaLM~\citep{Chowdhery:22short}, Llama~\citep{Touvron:23arxiv}, and Gemini~\citep{GeminiTeam:23arxiv,Reid:24arxiv}.

\section{Experiments}
\label{sec:Exp}

In this section, we present an empirical assessment of xLSTM, situating its performance within the context of established techniques, with particular attention to language modeling tasks. Section~\ref{sec:ExpSyntethic} explores the unique strengths of xLSTM through synthetic benchmarks. In Section~\ref{sec:ExpComparison}, we evaluate and contrast the perplexity on the validation set achieved by a variety of contemporary language modeling approaches, all of which have been trained on a dataset containing 15B tokens from SlimPajama~\citep{Soboleva:23}. Within the same experimental framework, we also undertake ablation studies to analyze the contributions of xLSTM components. Further, the scaling laws characteristic of the different models are examined following the methodologies outlined by \citet{Kaplan:20} and \citet{Brown:20short}.

Section~\ref{sec:ExpLanguage} expands the evaluation to a comprehensive language modeling experiment. Here, both xLSTM and the top-performing models from Section~\ref{sec:ExpComparison} are trained using a larger dataset of 300B tokens sourced from SlimPajama~\citep{Soboleva:23}. The subsequent analysis focuses initially on the capacity of these models to generalize to extended context lengths. This is followed by an assessment based on validation perplexity and results on downstream benchmarks~\citep{Sutawika:24}. Additionally, we measure model performance across 571 textual domains in the PALOMA language benchmark suite~\citep{Magnusson:23arxivshort}. Lastly, we revisit the scaling behaviors of each method under the increased data regime, encompassing a twentyfold enlargement of training examples.

\label{sec:xlstm_blocks_notation}
In all experiments, we denote xLSTM[$a$:$b$] as representing the configuration where the proportion of mLSTM-based to sLSTM-based xLSTM blocks is $a/b$. As an illustration, xLSTM[7:1] indicates that, out of a total of eight blocks, seven are constructed using mLSTM and one with sLSTM. When the overall block count is set to 48, this arrangement corresponds to 42 blocks of mLSTM type and 6 blocks of sLSTM type. Additionally, our experimental setup consistently incorporates pre up-projection blocks for mLSTM components and post up-projection blocks for sLSTM components.

\subsection{Synthetic Tasks and Long Range Arena}
\label{sec:ExpSyntethic}
Initially, we evaluate the capability of xLSTM’s exponential gating with memory mixing on formal languages as described by~\citet{Deletang:23}. Next, we investigate the advantages of xLSTM’s novel matrix memory using the Multi-Query Associative Recall task~\citep{Arora:23arxiv}. Lastly, xLSTM’s ability to handle extended sequences is examined on the Long Range Arena benchmark~\citep{Tay:21}.

\paragraph{Evaluation of xLSTM's Exponential Gating Combined with Memory Mixing.} We investigate the effectiveness of xLSTM's novel exponential gating mechanism coupled with memory mixing, designed to address state tracking challenges~\citep{Merrill:24, Merrill:23}. Building on the formal language benchmarks from \citet{Deletang:23}, we adapt and broaden these tasks to facilitate multi-length training, which allows for examination of the model's generalization across different sequence lengths. Full task specifications and expanded experimental outcomes are provided in Appendix~\ref{sec:appExpSynthetic-formalLang}. We benchmark xLSTM against various alternatives, including Transformers, State Space Models, and Recurrent Neural Networks. Each method's performance is quantified by accuracy on task-specific relevant tokens, normalized such that a value of 0 reflects chance and 1 denotes flawless performance. Comparisons are performed using 2-block configurations, including xLSTM[0:1] (corresponding to sLSTM only), xLSTM[1:0] (mLSTM only), xLSTM[1:1], Llama, Mamba, RWKV, Retention, Hyena, LSTM, as well as LSTM implemented within Transformer-style blocks (LSTM (Block)). Figure~\ref{fig:formal-main} presents the experimental findings. Methods lacking memory mixing—such as Transformers and State Space Models—are unable to solve state-dependent tasks, including regular grammars exemplified by the parity task. These outcomes are consistent with prior evidence that such models are inherently less expressive for these problems than RNNs~\citep{Merrill:24, Merrill:23, Deletang:23}.

\paragraph{Test of xLSTM's Memory Capacities on Associative Recall Tasks.} This study investigates the matrix-based memory enhancements of xLSTM by evaluating its memory capacity on the Multi-Query Associative Recall task~\citep{Arora:23arxiv}. Within each sequence, a set of key-value pairs is randomly selected from a large vocabulary and must be stored in memory for future retrieval. To raise the challenge of the standard setup, we increase the quantity of key-value pairs to as many as 256, and stretch the context window to 2048 tokens. These modifications allow a more extensive assessment of the memory profiles of several models. We include comparative evaluations of 2-block configurations from Llama, Mamba, RWKV-5, RWKV-6, xLSTM[1:1], and xLSTM[1:0]. Performance is quantified by the recall accuracy of these key-value pairs. Given that Transformer architectures such as Llama possess a memory scaling exponentially with their code dimension~\citep{Ramsauer:21}, they are considered the benchmark for this assessment. Figure~\ref{fig:mqar-main} presents the experimental outcomes. Remarkably, xLSTM[1:1] achieves the highest performance among all non-Transformer contenders, retaining this lead even for compact model sizes. Notably, the presence of the sLSTM block appears to not only maintain but actually bolster memory capacity, particularly evident at the uppermost difficulty level with 256 key-value pairs. Further findings, including extrapolation studies suggesting that xLSTM's memory advances facilitate generalization to sequence lengths beyond the training regime, are detailed in Appendix~\ref{sec:appExpSynthetic-mqar}.

\paragraph{Test of xLSTM's Long Context Capabilities on Long Range Arena.} In order to evaluate how well xLSTM manages extended sequences and substantial context windows, we benchmarked several approaches using the Long Range Arena~\citep{Tay:21}. Across all tasks, xLSTM consistently achieves robust results, indicating that its architecture is highly adept at addressing the challenges posed by long-context scenarios.

For more details, see Appendix~\ref{sec:appExpSynthetic-lra}.

\subsection{Method Comparison and Ablation Study}
\label{sec:ExpComparison}

In order to investigate the primary focus of our study—namely, the capabilities of our proposed LSTM variants when scaled for language modelling—we train xLSTMs, along with Transformers, State Space Models, and additional techniques, utilizing 15B tokens sourced from SlimPajama under consistent auto-regressive conditions. Model comparisons are conducted using the validation dataset, and comprehensive ablation analyses are carried out for the xLSTMs.

\paragraph{Benchmarking xLSTM Against Baseline Approaches.} To facilitate a comprehensive comparison, we train all models using a dataset of 15B tokens sourced from SlimPajama~\citep{Soboleva:23}. The performance metric used is the validation set perplexity. Our evaluation covers a variety of methodologies: xLSTM (the novel model presented here), GPT-3 (Transformer architecture)~\citep{Brown:20short}, Llama (Transformer architecture)~\citep{Touvron:23arxiv}, H3 (a State Space Model)~\citep{Fu:23}, Mamba (State Space Model)~\citep{Gu:24arxiv}, RWKV-4/5/6 (all Recurrent Neural Network variants)~\citep{Peng:23arxivshort, Peng:24arxivshort}, GLA (a linear Transformer variant)~\citep{Yang:23arxiv}, HGRN and HGRN2 (both RNN-based)~\citep{Qin:23, Qin:24arxiv}, as well as RetNet and Hyena (both classified as linear Transformers)~\citep{Sun:23arxiv, Poli:23}, and two variations, xLSTM[1:0] and xLSTM[7:1] (notation described in Section~\ref{sec:xlstm_blocks_notation}). 

Training was conducted using mixed precision; however, for RWKV-5, RWKV-6, GLA, and HGRN2, the mixed-precision implementation did not employ PyTorch's automatic mixed precision functionality (further procedural details appear in Appendix~\ref{sec:appGenTrainProc}). We group models into three broad categories: (a) Transformers, (b) State Space Models (SSMs), and (c) Recurrent Neural Networks (RNNs), together with linear Transformer approaches. Here, linear Transformers denote architectures replacing the standard attention mechanism with linear alternatives. 

Each model is configured to approximately align with GPT-3 at 350M parameters, specifically with an embedding dimension of 1024 and consisting of 24 residual blocks; it should be noted that GPT-3 uniquely utilizes tied embeddings for token input and output, which slightly reduces its parameter count. 

Per Table~\ref{tab:spaj15b_model_results}, xLSTM consistently achieves lower validation perplexity than all comparison methods. Further discussion and exhaustive experimental details are located in Appendix~\ref{sec:appExpComparison}. Additionally, Figure~\ref{fig:temp_sclaw15B} captures the scaling trends observed in this setup, illustrating that xLSTM is projected to scale effectively for models of larger size.

\paragraph{Ablation Studies.}
As presented in Table~\ref{tab:spaj15b_model_results} and Figure~\ref{fig:temp_sclaw15B}, xLSTM demonstrates superior language modeling performance when trained on 15B tokens from SlimPajama. To systematically analyze the individual contributions of the xLSTM enhancements, we gradually transform a standard LSTM model into the full xLSTM configuration. We begin by placing LSTM layers within pre-LayerNorm residual blocks. Next, we modify the model by introducing a post up-projection module. In the final step, we incorporate exponential gating and augment the architecture with matrix memory.

The results are shown in Table~\ref{tab:ablstudies} (top). 

The ablation experiments indicate that the observed performance gains stem from a combination of the exponential gating mechanism and the matrix memory. Moreover, considering the pivotal role that gating plays in RNNs and State Space Models, we systematically investigate various gating strategies. The results, presented in Table~\ref{tab:ablstudies} (bottom), demonstrate that allowing each gate to be learnable and conditioned on the input yields progressively improved outcomes. Further examination of individual backbone elements can be found in Appendix~\ref{sec:appAblDetails}.

\subsection{xLSTM as Large Language Model}
\label{sec:ExpLanguage}

This investigation concludes with comprehensive language modeling trials aimed at evaluating the suitability of xLSTM as a large language model. For this purpose, the scale of training is increased: models are trained on 300B tokens sourced from SlimPajama. This token count aligns with training regimens utilized in other recent works, including Mamba~\citep{Gu:24arxiv} and Griffin~\citep{De:24arxiv}.

A comparative analysis is performed between xLSTM and three established baselines: RWKV-4, Llama, and Mamba, each representing a distinct architecture discussed in Section~\ref{sec:ExpComparison}. RWKV-4 is chosen as the RNN exemplar since appropriate precision configurations for RWKV-5, RWKV-6, and HGRN2 were finalized only after commencement of the 300B token training, as detailed in Appendix~\ref{sec:appGenTrainProc}. 

Models are trained across a range of capacities (125M, 350M, 760M, 1.3B). Each model undergoes rigorous testing for its ability to extrapolate to longer sequence lengths, and their performance is validated using the designated evaluation set. Further, models are analyzed on downstream task benchmarks, with language modeling proficiency measured on 571 distinct text domains comprising the PALOMA benchmark. The examination concludes with an analysis of scaling law behavior across the suite of models.

\paragraph{Sequence Length Extrapolation.}
\label{sec:spaj300B_ctx_extrapolation}
We begin by evaluating the ability of large-scale, 1.3B parameter models—including xLSTM, RWKV-4, Llama, and Mamba—to generalize to longer sequence lengths. Each model is trained using a maximum context length of 2048, after which inference is conducted on increasingly longer context sizes, reaching up to 16384. The outcomes of these experiments are presented in Figure~\ref{fig:spaj300B_seq_extrapolation}. Notably, xLSTM consistently achieves lower perplexity at extended context lengths compared to competing approaches.

\paragraph{Validation Perplexity and Downstream Tasks.}
\label{sec:spaj_downstream_eval}
Next, for every examined model scale, we assess xLSTM, RWKV-4, Llama, and Mamba by measuring their next token prediction performance using the SlimPajama validation split, as well as their abilities on downstream benchmarks that focus on aspects of common sense reasoning.

The third column in Table~\ref{tab:lmeval} displays the validation set perplexities achieved by various approaches. Across all model capacities, xLSTM[1:0] and xLSTM[7:1] outperform other configurations in terms of validation perplexity. The subsequent columns in Table~\ref{tab:lmeval} summarize the results on downstream evaluation benchmarks. Generally, xLSTM consistently attains superior performance across almost all downstream tasks and model scales, with the exception of certain cases on the ARC dataset, where Mamba shows better results. Further information is provided in Appendix~\ref{sec:appExpLanguage}.

\paragraph{Performance on PALOMA Language Tasks.} Third, we evaluate the capabilities of xLSTM, RWKV-4, Llama, and Mamba across all model scales by examining their accuracy in next token prediction on PALOMA language tasks~\citep{Magnusson:23arxivshort}. For this analysis, we use perplexity as the primary metric, measuring performance on 571 distinct textual domains that range from outlets like nytimes.com to communities such as r/depression on Reddit. Table~\ref{tab:ppleval} summarizes the results, dividing them between general language modeling benchmarks (first seven columns) and fine-grained, domain-specific assessments (last five columns). In this setting, xLSTM[1:0] consistently outperforms xLSTM[7:1].

Comparing against the other models, xLSTM[1:0] exhibits lower perplexity than Mamba on 568 of 571 domains (99.5\%), outperforms Llama on 486 of 571 (85.1\%), and surpasses RWKV-4 in 570 of 571 domains (99.8\%). Additional information can be found in Appendix~\ref{sec:appExpLanguage}.

\paragraph{Scaling Laws.} Fourth, we examine the power-law scaling properties that facilitate projecting model performance at larger scales~\citep{Kaplan:20,Brown:20short}. The scaling trends are illustrated in Figure~\ref{fig:temp_sclaw300B}. Although all evaluated models display comparable scaling tendencies, they differ in their vertical intercepts. Among them, RWKV-4 achieves the lowest performance, with Llama and Mamba performing incrementally better, respectively. xLSTM surpasses Mamba by a margin similar to that which Mamba holds over Llama. These observed scaling patterns suggest that as model size increases, xLSTM is likely to sustain its advantage relative to both Transformers and State-Space models.

\textbf{Generation Times and Maximal Throughput.} We report the time required for text generation in Figure~\ref{fig:inference_time_generation} and present the maximal throughput results in Figure~\ref{fig:inference_time_decoding_throughput} (left) for the various xLSTM variants at the 1.3B parameter scale. These are evaluated in comparison to similarly sized versions of Mamba, Llama, and RWKV implemented via HuggingFace, including Llama equipped with a static key-value cache. During experimentation, attempts at complete cache compilation for the Transformer Model, as well as compiling Mamba with \texttt{torch.compile}, were unsuccessful. All generation experiments employ a batch size of 1 with a pre-fill length of 16, a configuration that most benefits the Transformer architecture. 
The results in Figure~\ref{fig:inference_time_generation} illustrate that xLSTM, along with the recurrent Mamba and RWKV-4 models, exhibits linear scaling properties, in contrast to the quadratic scaling observed with Llama. In our evaluation of decoding throughput, multiple batch sizes and pre-fill settings were investigated for the Llama model.
Figure~\ref{fig:inference_time_decoding_throughput} (right) demonstrates that xLSTM is able to handle substantially larger batch sizes than Llama, a consequence of its constant memory footprint, thereby attaining the best throughput among the models considered.

\section{Limitations}

(i) Unlike mLSTM, the structure of memory mixing in sLSTM inherently restricts the use of parallel operations, making efficient parallel execution infeasible. Despite this, we have implemented a specialized CUDA kernel for sLSTM that achieves performance within a factor of two compared to our parallel implementation of mLSTM. (ii) Currently, our CUDA kernels for mLSTM are not fully optimized; as a result, mLSTM runs approximately four times slower than either FlashAttention or Mamba’s scan. With further optimization, similar to the strategies used in FlashAttention, faster mLSTM CUDA kernels should be attainable. (iii) The memory in mLSTM, represented as $d \times d$ matrices, introduces a substantial computational burden. However, as the memory read and write procedures are parameter-free and are implemented with conventional matrix operations, parallelization mitigates wall clock overhead, so the impact of the memory’s complexity remains limited. (iv) Careful consideration is required when initializing the forget gates. (v) Because the matrix memory design in mLSTM is decoupled from sequence length, scaling up the sequence length for longer contexts could potentially saturate the memory, though we observed no adverse effects for context sizes as large as 16k tokens (refer to Section~\ref{sec:spaj300B_ctx_extrapolation}). (vi) Given the prohibitive compute demands of running large language model experiments, comprehensive tuning of both architecture and hyperparameters—especially for xLSTM at larger scales—has not been conducted. We expect that more extensive optimization will be necessary for xLSTM to achieve optimal performance.

\section{Conclusion}
We have provided a partial answer to our foundational query: To what extent can language modeling be advanced by scaling LSTM architectures to billions of parameters? Our findings demonstrate that these models can, at a minimum, match the performance of prevailing methods such as Transformers and State Space Models. Through the introduction of xLSTM—which incorporates exponential gating, memory mixing, and a novel memory architecture—we observe superior results in language modeling tasks when compared against leading approaches like Transformers and State Space Models. The scaling behavior of these networks suggests that larger xLSTM models have the potential to rival current Large Language Models that rely on Transformer-based architectures. Beyond language modeling, the capabilities of xLSTM present promising possibilities for broader applications, including Reinforcement Learning, Time Series Prediction, and the simulation of physical systems.

\end{document}